% =====================================================================
% CAPÍTULO 6: DISCUSSÃO (piso de rodadas)
% Interpreta os resultados do Capítulo 5, sem repeti-los.
% Seções: 6.1 par de mensagens; 6.2 origem do piso; 6.3 o piso e os
% detectores; 6.4 literatura; 6.5 significado prático;
% 6.6 limitações reveladas pelos resultados.
% =====================================================================

\chapter{Discussão}
\label{cap:discussao}

Este capítulo interpreta os resultados do Capítulo~\ref{cap:resultados}. A Seção~\ref{sec:disc-par} explica por que um par de mensagens revela a cifra com mais rodadas que um criptograma isolado, e as seções seguintes tratam de onde vem o piso, da relação entre o piso e os detectores, da comparação com a literatura, do que o piso significa na prática e das limitações que os resultados revelaram.

\section{Por que o par de mensagens revela a cifra}
\label{sec:disc-par}

Nos três algoritmos, cada bloco cifrado é o bloco de texto em claro somado a uma máscara, que depende da chave, do nonce e das rodadas que o estado atravessou até ali, $C = P \oplus M(K, n)$~\cite{nist_sp800232_2025,Banik2021,Beierle2021SparkleSub}. Uma cifra segura produz máscaras que não podem ser separadas de bytes aleatórios sem a chave~\cite{Rogaway2004Nonce,KatzLindell2014IMC}, e é essa propriedade que as rodadas garantem. Com poucas rodadas, a máscara ainda carrega estrutura, mas o adversário passivo não a observa diretamente~(Seção~\ref{sec:met-adversario}), ele vê o criptograma, em que a máscara aparece somada a um texto que ele não conhece e calculada com uma chave que ele também não conhece.

No par, as duas mensagens são cifradas com a mesma chave e com nonces consecutivos, que, nos pares usados, diferem em um único bit, e o XOR dos dois criptogramas é
\begin{equation}
C_1 \oplus C_2 = (P_1 \oplus P_2) \oplus \big(M(K, n) \oplus M(K, n+1)\big).
\end{equation}
O primeiro termo é a combinação dos dois textos em claro, e o segundo, a diferença entre as duas máscaras. Com o número de rodadas da especificação, a diferença entre os nonces se espalha por todo o estado antes de chegar ao texto, e as duas máscaras ficam tão diferentes entre si quanto duas máscaras sorteadas. Com poucas rodadas, essa diferença ainda não se espalhou, as duas máscaras coincidem em boa parte das posições, e a parte em que coincidem se cancela no XOR, junto com a dependência da chave que ela carrega. O que sobra é a diferença que as rodadas não espalharam~\cite{BihamShamir1991}, que nasce de uma diferença fixa no nonce e, por isso, depende menos da chave do que cada máscara isolada. Isso explicaria por que um classificador treinado com chaves próprias reconhece o par em dispositivos cujas chaves ele nunca viu. O teste com nonces a distância $k$ é coerente com essa explicação, já que, quando os dois nonces diferem em vários bits, o F1 caiu e, no Ascon-AEAD128 e no GIFT-COFB, o piso perdeu uma rodada (Seção~\ref{sec:res-contador}).

O primeiro termo, $P_1 \oplus P_2$, também não é aleatório. Dois trechos de texto somados bit a bit mantêm regularidades dos dois, como um bit mais significativo fixo em cada byte, que, em texto ASCII, por exemplo, é zero e continua zero no XOR, e é esse tipo de estrutura que permite recuperar os textos quando duas mensagens são cifradas com a mesma máscara~\cite{Mason2006TwoTimePad}. O piso foi o mesmo com o corpus original, que inclui marcação HTML e livros em outros idiomas, e não diminuiu quando 20\% das mensagens foram trocadas por imagens, com uma única célula detectada no Ascon-AEAD128 com quatro rodadas, com F1 de 0,512 (Seção~\ref{sec:res-robustez}). Aqui, as máscaras não são iguais, só parecidas, e o classificador não recupera os textos, mas a soma dos dois termos continua distinguível de bytes aleatórios enquanto a diferença entre as máscaras for pequena. O braço aleatório confirma o papel do texto, já que, com texto em claro aleatório, $P_1 \oplus P_2$ é uniforme e esconde qualquer diferença entre as máscaras na parte cifrada. Nesse braço houve detecções isoladas, mas a proporção de células com p-valor abaixo de 0,05 ficou dentro do esperado ao acaso, e nenhum dos três algoritmos teve piso em nenhum número de rodadas (Seção~\ref{sec:res-calibracao}).

Com um criptograma isolado, não há uma segunda mensagem da mesma chave com que compará-lo, e a máscara não se cancela. A sua máscara depende da chave, e cada dispositivo de teste tem uma chave que o classificador nunca viu, de modo que aprender como é a máscara de uma chave não ajuda a reconhecer a de outra. O par anula a chave sem conhecê-la, e o criptograma isolado só revela a cifra quando a própria máscara ainda carrega estrutura que não depende da chave, o que acontece com uma única rodada (Seção~\ref{sec:res-isolado-80b}). Com uma rodada, parte da máscara vem do valor inicial, das constantes e do nonce, que, com o contador começando do zero, é quase o mesmo em todas as mensagens, e a chave ainda não se misturou de forma não linear a todas as posições. O caso extremo é o Schwaemm256-128 com um passo, separado dos dados aleatórios com F1 de 0,999 por mensagem. Com duas rodadas, a chave já alcança toda a máscara, e o criptograma isolado deixa de ser detectado no Ascon-AEAD128 e no GIFT-COFB e passa a ser detectado só de forma marginal no Schwaemm256-128. Essa é uma explicação possível para os resultados, mas não foi medida diretamente, o que exigiria comparar as máscaras produzidas por chaves diferentes com o mesmo nonce.

Com criptogramas de 64~KB descritos pelos atributos estatísticos, nem essa primeira rodada foi detectada (Seção~\ref{sec:res-isolado-64k}). Uma explicação possível é que a parte mal misturada fica no início da mensagem, já que o estado acumula rodadas a cada bloco, e os atributos, calculados sobre os 65\,552 bytes do criptograma, misturam essa região com milhares de blocos indistinguíveis de bytes aleatórios. O Schwaemm256-128 com um passo é compatível com essa hipótese, com F1 de 0,999 nos 80 bytes descritos pelos bits crus e no acaso nos 64~KB descritos pelos atributos. Como, de um experimento para o outro, mudaram ao mesmo tempo o tamanho da mensagem e a representação, a diluição é uma hipótese, e não uma causa demonstrada, mas foi ela que motivou concentrar a análise nos primeiros bytes da mensagem, bit a bit.

A comparação entre as duas formas de apresentar o par confirma que a fonte do sinal é observar duas mensagens da mesma chave, e não o XOR (Seção~\ref{sec:res-piso}). Com os dois criptogramas lado a lado, como na entrada do distinguidor de Gohr~\cite{gohr_2019}, o piso já é o mesmo do XOR nos três algoritmos, mas a rodada do piso é detectada em menos células e com sinal mais fraco, com F1 de 0,554 por par no Ascon-AEAD128 com três rodadas, contra 0,669 com o XOR. Com o XOR, como em Shen et al.~\cite{Shen2024ND}, o classificador recebe pronta a comparação bit a bit entre as duas mensagens, que, com os criptogramas lado a lado, ele precisa aprender, combinando cada bit de uma mensagem com o bit correspondente da outra. O XOR, portanto, não cria o sinal, ele o torna mais fácil de aprender, e o piso fica onde a diferença entre as máscaras deixa de ser visível, qualquer que seja a forma de apresentar o par. A seção seguinte trata de onde fica esse ponto e por quê.

\section{De onde vem o piso}
\label{sec:disc-origem}

A hipótese mais forte é que o sinal venha da rodada de inicialização. No Ascon-AEAD128 e no Schwaemm256-128, reduzir só a inicialização e a finalização reproduziu o piso e a força do sinal do estudo principal, com F1 de 0,667 por par no Ascon-AEAD128 com três rodadas, contra 0,669, e o mesmo aconteceu sem a tag, que é a única parte do criptograma produzida pela finalização, enquanto reduzir só as rodadas entre blocos não produziu detecção em nenhuma das 72 células do Ascon-AEAD128 nem das 63 do Schwaemm256-128, nem com uma única rodada entre blocos (Seção~\ref{sec:res-fase}). Essa hipótese é coerente com o que a Seção~\ref{sec:disc-par} mostrou. O par só tem sinal enquanto as duas máscaras forem parecidas, e elas só são parecidas enquanto a diferença do nonce não tiver se espalhado. Com a inicialização completa, essa diferença se espalha antes do primeiro bloco de texto, as duas máscaras já saem independentes, e não resta diferença para as rodadas entre blocos preservarem, por mais fracas que sejam.

A hipótese, portanto, é que o piso dependa menos do número de rodadas de cada chamada que de quantas rodadas a diferença do nonce atravessa antes de encontrar o texto. O GIFT-COFB ilustra essa hipótese. Sem dado associado, a cifra de bloco é chamada duas vezes antes do primeiro bloco de texto, uma para o nonce e outra para o bloco de preenchimento do dado associado, e o piso de três rodadas corresponde, na prática, a seis rodadas da cifra de bloco no caminho da diferença.

A posição do contador no nonce muda o piso porque muda o ponto em que a diferença entre os nonces entra no estado. Com o contador nos primeiros bytes do nonce, o Ascon-AEAD128 com três rodadas ficou mais forte, com F1 de 0,893 por par, contra 0,669, e o GIFT-COFB e o Schwaemm256-128 perderam uma rodada (Seção~\ref{sec:res-contador}). No Ascon-AEAD128, o nonce ocupa duas palavras de 64 bits do estado, e a posição do bit que muda define em qual coluna de S-boxes a diferença começa e quais rotações da camada linear ela encontra nas primeiras rodadas. Nos três algoritmos, a diferença percorre caminhos diferentes conforme a posição, e alguns se espalham mais devagar que outros. Essa explicação é coerente com os resultados, mas o caminho da diferença dentro do estado não foi rastreado.

Os pisos de 3, 3 e 2 também não ordenam os algoritmos por segurança, porque uma rodada não é a mesma unidade nos três. Uma rodada do Ascon-AEAD128 aplica uma camada de S-boxes de 5 bits e uma camada linear sobre um estado de 320 bits, uma rodada do GIFT-128 aplica uma camada de S-boxes de 4 bits e uma permutação de bits sobre 128 bits, e um passo do SPARKLE aplica a cada par de palavras uma caixa ARX de quatro rodadas, a Alzette, seguida de uma camada linear~\cite{Beierle2020Sparkle}. O caso do GIFT-COFB e o efeito da posição do contador mostram, ainda, que o mesmo número de rodadas pode corresponder a caminhos de comprimentos diferentes até o texto. O que cada piso permite comparar é o algoritmo reduzido com a sua própria especificação.

O piso fica onde o sinal da cifra cruza o limite do instrumento. A cifra define quanto o sinal cai a cada rodada, e o instrumento define em que ponto dessa queda está o limiar. Um adversário com mais dados, ou com um detector que saiba onde procurar, pode ir além. A seção seguinte discute a relação entre o piso e os detectores.


\section{O piso e os detectores}
\label{sec:disc-detectores}

Os detectores avaliados vão do mais simples, a regressão logística, que pondera cada bit separadamente, a redes neurais treinadas com dez vezes os dados do estudo principal, e todos chegaram à mesma rodada. No Ascon-AEAD128 com três rodadas, a regressão logística ficou entre 0,661 e 0,665 por par, contra 0,669 do melhor classificador, e, no Schwaemm256-128 com duas rodadas, entre 0,946 e 0,947, contra 0,998 (Seções~\ref{sec:res-xor} e~\ref{sec:res-dados}). A busca de hiperparâmetros, as redes neurais e dez vezes mais dados não subiram o piso de nenhum algoritmo com os detectores testados (Seção~\ref{sec:res-robustez}). Entre os detectores avaliados, portanto, o piso não dependeu de um modelo específico, e o fato de um detector simples chegar ao mesmo ponto que os mais expressivos indica que ele reflete a cifra vista por esse adversário, e não a escolha do detector.

Os modelos mais expressivos fizeram diferença na força do sinal, e não no piso, e só no GIFT-COFB. Com três rodadas, as árvores chegaram a 0,598 por par, contra 0,524 a 0,533 da regressão logística, e a MLP, a 0,553, contra 0,533 com a mesma quantidade de dados. Árvores e redes representam combinações de bits, enquanto a regressão logística pondera cada bit separadamente, o que sugere que, no GIFT-COFB, parte do sinal esteja em combinações de bits. No Ascon-AEAD128, as redes ficaram abaixo da regressão logística, com 0,621 e 0,642, contra 0,665.

Esse comportamento é coerente com a literatura de distinguidores neurais. Gohr~\cite{gohr_2019} mostrou que uma rede neural supera o distinguidor puramente diferencial no Speck32/64 com rodadas reduzidas, e Benamira et al.~\cite{benamira_2021} mostraram depois que é possível reproduzir o desempenho dessa rede com um distinguidor construído a partir de características criptanalíticas explícitas. Gohr, Leander e Neumann~\cite{GohrLeanderNeumann2022} mostraram, ainda, que os ganhos atribuídos a distinguidores neurais sobre vários pares em geral desaparecem diante de uma combinação clássica das respostas de pares independentes. Perto do piso, a vantagem de um modelo expressivo sobre um detector simples foi pequena, e só apareceu no GIFT-COFB.

% >>> OCULTO TEMPORARIAMENTE (envio ao orientador só até a 6.3).
% Para desocultar, apagar esta linha \iffalse e o \fi no fim do arquivo.

\section{Comparação com a literatura}
\label{sec:disc-literatura}

Os trabalhos que medem até quantas rodadas os finalistas resistem chegam, em geral, mais longe que este, e a razão está no adversário. A Tabela~\ref{tab:disc-literatura} resume, para cada trabalho, o objeto analisado, o que o adversário controla, o que ele observa e até onde a detecção chegou.

A maior parte da literatura escolhe a diferença de entrada. Baksi et al.~\cite{Baksi2021} construíram distinguidores diferenciais com aprendizado de máquina para a permutação do Ascon com três rodadas, e Shen et al.~\cite{Shen2024ND} chegaram a quatro rodadas na mesma permutação e a sete das 40 rodadas na cifra de bloco GIFT-128, com uma diferença de um bit escolhida na entrada da primitiva isolada, sem chave no caso do Ascon e sem nonce nem encadeamento no caso do GIFT-128. Rajan et al.~\cite{Rajan2022} aplicaram a mesma abordagem ao GIFT-COFB, o trabalho mais próximo deste no objeto, e distinguiram de dados aleatórios versões com 2 a 6 rodadas, também com diferenças de entrada escolhidas. Na criptanálise clássica do Ascon, Dobraunig et al.~\cite{Dobraunig2021AsconSecurity} usaram uma característica diferencial-linear de quatro rodadas da inicialização numa recuperação de chave, Civek e Tezcan~\cite{CivekTezcan2022} chegaram a cinco rodadas da permutação, e Rohit et al.~\cite{Rohit2021Ascon} chegaram a sete rodadas da inicialização, com uma recuperação de chave que respeita o limite de dados da especificação, mas custa cerca de $2^{123}$ chamadas da permutação. Em todos esses trabalhos, o adversário escolhe as entradas ou os nonces e observa a saída da primitiva ou a keystream, sem um texto desconhecido somado a ela.

O trabalho mais próximo deste na pergunta é o de Bellini, Huang e Rachidi~\cite{BelliniHuang2022}, que aplicaram testes estatísticos de aleatoriedade, entre eles os do NIST~\cite{Rukhin2010NIST80022}, às versões com rodadas reduzidas das primitivas dos finalistas, sob várias estratégias de geração de dados. Na estratégia com entradas aleatórias, a única em que o analista não controla a entrada, os testes detectaram apenas a primeira rodada na maior parte das primitivas, entre elas a permutação do Ascon, as permutações SPARKLE e o GIFT-128, enquanto, com entradas escolhidas, chegaram a oito rodadas no GIFT-128. Com um criptograma isolado de 80 bytes, este trabalho também só detectou a primeira rodada no Ascon-AEAD128 e no GIFT-COFB, o que é coerente com esse resultado.

O par de mensagens fica entre os dois regimes. O adversário não escolhe nada, mas o contador do nonce lhe entrega, sem custo, uma diferença de entrada conhecida e pequena entre duas mensagens consecutivas, e é isso que leva a detecção da primeira para a terceira rodada no Ascon-AEAD128 e no GIFT-COFB.

Nos estudos de identificação de algoritmos criptográficos por aprendizado de máquina reunidos na revisão sistemática (Capítulo~\ref{cap:trabalhos-relacionados}), nenhum avalia esquemas de criptografia autenticada nem versões com rodadas reduzidas. O resultado nulo com criptogramas de 64~KB, mesmo com uma rodada, indica que o desenho usado nesses estudos, atributos estatísticos calculados sobre o criptograma inteiro, não alcançou o sinal, que, como hipótese, estaria no início da mensagem (Seção~\ref{sec:disc-par}), e este trabalho acrescenta a essa linha a pergunta de quantas rodadas são necessárias para que esse nulo apareça.

Nenhum desses resultados é contradito pelos deste trabalho. Os números da literatura ficam acima do piso medido aqui, como se espera de adversários que controlam as entradas, e o piso com o par fica acima do obtido por Bellini et al. com entradas aleatórias, como se espera de um adversário que observa duas mensagens com nonces relacionados.

\begin{table}[htb]
\centering
\caption{Comparação com trabalhos que medem o número de rodadas detectável nos finalistas.}
\label{tab:disc-literatura}
\scriptsize
\setlength{\tabcolsep}{3pt}
\begin{tabular}{p{2.6cm}p{2.6cm}p{2.4cm}p{2.6cm}p{2.4cm}}
\toprule
Trabalho & Objeto & O adversário controla & O adversário observa & Rodadas detectadas \\
\midrule
Baksi et al.~\cite{Baksi2021} & permutação do Ascon & diferença de entrada & saída da permutação & 3 de 12 \\
Shen et al.~\cite{Shen2024ND} & permutação do Ascon e GIFT-128 & diferença de entrada & saída da primitiva & 4 de 12 e 7 de 40 \\
Rajan et al.~\cite{Rajan2022} & GIFT-COFB & diferença de entrada & diferença de saída do cifrador & 6 de 40 \\
Dobraunig et al.~\cite{Dobraunig2021AsconSecurity} & inicialização do Ascon & diferença no nonce & keystream & 4 de 12 \\
Civek e Tezcan~\cite{CivekTezcan2022} & permutação do Ascon & diferença de entrada & saída da permutação & 5 de 12 \\
Rohit et al.~\cite{Rohit2021Ascon} & inicialização do Ascon & nonces & keystream & 7 de 12 \\
Bellini et al.~\cite{BelliniHuang2022} & primitivas dos finalistas & nada (entrada aleatória) & saída da primitiva & 1 \\
Este trabalho, criptograma isolado & Ascon-AEAD128, GIFT-COFB e Schwaemm256-128 & nada & um criptograma, texto desconhecido & 1, 1 e 2 \\
Este trabalho, par de mensagens & Ascon-AEAD128, GIFT-COFB e Schwaemm256-128 & nada & dois criptogramas consecutivos, texto desconhecido & 3, 3 e 2 \\
\bottomrule
\end{tabular}
\end{table}

\section{O que o piso significa na prática}
\label{sec:disc-pratica}

O piso permite uma afirmação num único sentido. Uma versão com o número de rodadas do piso, ou menos, está descartada, porque um adversário passivo, que não conhece a chave nem o texto, já a distingue de dados aleatórios, o que, neste trabalho, vale para até três rodadas no Ascon-AEAD128, até três no GIFT-COFB, que correspondem a seis rodadas da cifra de bloco antes do primeiro bloco de texto, e até dois passos no Schwaemm256-128. Acima do piso, o trabalho não afirma nada sobre a segurança, já que o piso é um limite inferior. O piso indica, portanto, até onde um algoritmo não pode ser aliviado, e não quanto ele pode.

Na rodada do piso, explorar a distinção custa pouco ao adversário. Com bolsas de um par, em que cada decisão usa duas mensagens consecutivas, o F1 sobre o conjunto de teste já fica acima do acaso, com um classificador treinado com 1 a 2\,000 pares, mas cada decisão isolada ainda erra com frequência, e, para decidir com confiança sobre um dispositivo, são necessárias bolsas de cerca de 100 pares, ou 200 mensagens, com 10\,000 a 30\,000 pares de treinamento no Ascon-AEAD128 e no GIFT-COFB (Seção~\ref{sec:res-piso-dados}). Num dispositivo que envia uma mensagem por segundo, 200 mensagens correspondem a pouco mais de três minutos de tráfego. O treinamento não depende do dispositivo-alvo, porque o adversário pode gerar os pares de treino com uma implementação da versão reduzida e chaves próprias. Na rodada do piso, o risco não é, portanto, apenas teórico.

Diante desse adversário, as especificações têm margem grande, com 12 rodadas de inicialização no Ascon-AEAD128 contra um piso de três, 40 rodadas por chamada no GIFT-COFB contra três, e 11 passos no Schwaemm256-128 contra dois. Essa margem não autoriza aliviar os algoritmos até perto do piso, porque adversários mais fortes chegam mais longe, como mostra a Seção~\ref{sec:disc-literatura}, com ataques de sete rodadas ao Ascon. A margem que importa para o projeto de uma versão aliviada é a do adversário mais forte que se queira considerar, e o piso deste trabalho é o ponto dessa escala que corresponde a um adversário que apenas observa o tráfego. Ele é a referência adequada quando esse é, de fato, o modelo de ameaça do dispositivo.

As condições de uso também mudam o piso. Entre as posições testadas, o contador nos últimos bytes do nonce foi a mais favorável ao adversário, já que, com o contador nos primeiros bytes, o piso não aumentou em nenhum algoritmo e diminuiu no GIFT-COFB e no Schwaemm256-128, e pares cujos nonces diferem em vários bits deram um F1 menor e, no Ascon-AEAD128 e no GIFT-COFB, um piso uma rodada menor (Seção~\ref{sec:res-contador}). O texto em claro também pesa, porque o sinal só existe enquanto ele tem padrão, e uma fonte com menos padrão, como dados comprimidos, pode dar um piso menor, enquanto cabeçalhos fixos de protocolo podem dar um piso maior (Seção~\ref{sec:met-escopo}).

\section{Limitações reveladas pelos resultados}
\label{sec:disc-limitacoes}

As limitações que decorrem do desenho do experimento, como a fonte de texto, o tamanho das mensagens, a ausência de dado associado e o contador começando do zero, foram apresentadas na Seção~\ref{sec:met-escopo}. Esta seção trata das que só apareceram nos resultados e indica, para cada uma, em que direção ela pode mover os números e se isso foi verificado.

Dois pisos dependem de uma única célula. No criptograma isolado de 80 bytes, o piso de dois passos do Schwaemm256-128 vem de uma célula com bolsas de 100 mensagens e F1 de 0,621, enquanto, com bolsa de uma mensagem, a mesma rodada teve F1 de 0,481 e AUC de 0,476, abaixo do acaso (Seção~\ref{sec:res-isolado-80b}). Quando a correção de Benjamini-Hochberg é refeita sem as células com F1 acima de 0,9, que são detecções certas e afrouxam o limite das demais, essa célula deixa de passar, e o piso isolado do Schwaemm256-128 cai para um passo, embora o teste de permutação por chave a mantenha, com p-valor de 0,001. O piso isolado de uma rodada do GIFT-COFB também vem de uma única célula, mas com p-valor da ordem de $10^{-6}$, e se mantém nessa verificação. No piso de dados, a detecção do Ascon-AEAD128 com um único par de treinamento teve F1 de 0,507, no limite do que o critério aceita (Seção~\ref{sec:res-piso-dados}). Se algum desses resultados for falso, o valor correto é menor. Os pisos do estudo principal, detectados em nove, oito e nove das nove células, não mudaram com essa verificação.

Os modelos foram treinados com uma única semente, e a separação entre as 240 chaves de treinamento e as 60 de teste foi sorteada uma única vez. Os intervalos de confiança, obtidos por reamostragem das chaves de teste, cobrem a variação do conjunto de teste, mas não a do treinamento nem a da separação. Células detectadas em uma ou duas das nove combinações de classificador e bolsa podem mudar com outra semente ou outra separação, enquanto rodadas detectadas nas nove células dificilmente mudam. Essa variação não foi medida, o que exigiria repetir o estudo com várias sementes e separações e reportar a frequência com que cada célula é detectada.

O controle com texto aleatório também tem um limite. A tag é calculada na finalização, que também é reduzida, e pode ter estrutura mesmo com texto aleatório, de modo que as dez detecções no braço aleatório, todas com a tag incluída, não podem ser tratadas como falsos positivos sem uma nova análise sem a tag, que não foi feita para essas variações. No estudo principal, a análise sem a tag mostrou um falso positivo em 180 células, e o teste A/A, uma taxa de p-valores abaixo de 0,05 próxima da esperada, sem nenhuma detecção depois da correção (Seção~\ref{sec:res-calibracao}).

Nem toda diferença de F1 pode ser interpretada. Com bolsas de 100 pares, cada célula tem apenas 120 decisões, 60 por classe, e o F1 oscila cerca de 0,09 para cima ou para baixo só por acaso, tanto que as rodadas sem detecção chegam a 0,571 e 0,574 no GIFT-COFB, com 20 e 40 rodadas, e a 0,582 no Ascon-AEAD128 com quatro rodadas. Diferenças pequenas nessa coluna, como entre 0,899 e 0,933, não indicam que um algoritmo seja mais fraco que outro. Algumas comparações também juntam números obtidos com conjuntos de teste de tamanhos diferentes, de 60, 600 e 1\,000 dispositivos, cuja precisão não é a mesma.

A conclusão de que o piso não depende do detector vale para os detectores testados. A CNN 1D descarta a posição de cada byte, o que limita o seu uso quando o sinal está em posições fixas, e, no GIFT-COFB com três rodadas, ela passou a prever quase sempre a mesma classe, o que indica falha de treinamento, e não ausência de sinal (Seção~\ref{sec:res-modelos}). As redes também foram treinadas sem parada antecipada, com um número fixo de épocas.

Parte das explicações da discussão foi inferida dos resultados, e não medida diretamente. A origem do sinal na inicialização, já que nenhum estudo reduziu só a inicialização (Seção~\ref{sec:disc-origem}), a estrutura da máscara que permite detectar um criptograma isolado com uma rodada (Seção~\ref{sec:disc-par}), o caminho da diferença conforme a posição do contador (Seção~\ref{sec:disc-origem}) e a diluição do sinal nos criptogramas de 64~KB são explicações possíveis, e cada uma pode ser confirmada ou refutada por uma medição simples, como comparar as máscaras produzidas por chaves diferentes com o mesmo nonce, acompanhar a diferença rodada a rodada dentro do estado ou medir o sinal bloco a bloco ao longo da mensagem.

Por fim, o critério de detecção foi definido neste trabalho. O teto da faixa de controle e a exigência de que todas as rodadas abaixo do piso sejam detectadas o tornam mais rigoroso, e, como o p-valor é tomado como 1 quando o F1 fica abaixo de 0,50, o teste é, na prática, unilateral (Seção~\ref{sec:met-criterio}). Um critério menos rigoroso poderia apontar pisos maiores, embora, no estudo principal, nenhuma célula acima do piso tenha passado pela correção de Benjamini-Hochberg, de modo que, ali, os pisos seriam os mesmos sem o teto e sem a regra de monotonicidade, e a verificação descrita acima mostra que um critério mais rigoroso mantém os pisos do estudo principal e altera apenas o piso isolado do Schwaemm256-128.


% <<< FIM DO TRECHO OCULTO
